\section{Data and methodology}

\subsection{Dataset and index membership}

The dataset contains three cross-sections of FTSE/JSE Top 40 instruments observed or reconstructed for 30 June 2025, 31 December 2025 and 30 June 2026. Each cross-section contains 41 listed instruments representing 40 company groups. The difference arises because Investec Limited (INL) and Investec plc (INP) are separate listed securities of one dual-listed company. Both instruments were retained because the analysis compares listed security profiles. Table~\ref{tab:1} summarises the population at each decision date.


\begin{table}[htbp]
\centering\caption{Date-specific population reconciliation.}\label{tab:1}
\footnotesize
\begin{tabularx}{\linewidth}{@{}>{\hsize=0.880\hsize\linewidth=\hsize\raggedright\arraybackslash}X >{\hsize=0.520\hsize\linewidth=\hsize\raggedright\arraybackslash}X >{\hsize=0.520\hsize\linewidth=\hsize\raggedright\arraybackslash}X >{\hsize=2.080\hsize\linewidth=\hsize\raggedright\arraybackslash}X@{}}
\toprule
Decision date & Instruments & Groups & Change since preceding date \\
\midrule
June 2025 & 41 & 40 & Initial observed population \\
December 2025 & 41 & 40 & GLN admitted; APN removed \\
June 2026 & 41 & 40 & PAN admitted; MNP removed \\
\bottomrule
\end{tabularx}
\end{table}


Index membership was reconciled against the effective-dated quarterly reviews. Northam Platinum replaced Glencore before the June 2025 observation. Glencore returned in September 2025, replacing Aspen, while the December review recorded no Top 40 change. Pan African Resources replaced Mondi in March 2026, and no change was recorded in June 2026. These changes are important because an instrument entering or leaving the index should not be treated as an ordinary change in financial peer group. Figure~\ref{fig:population} shows the date-specific populations and the instruments available for comparison.



\begin{figure}[htbp]
\centering
\diagram{population}
\caption{Index populations and continuity. The flow population has 39 complete histories; every date-specific fit retains all 41 instruments.}\label{fig:population}
\end{figure}


The market variables use trailing annual windows, as shown in Figure~\ref{fig:windows}. Adjacent windows overlap by six months.


\begin{figure}[htbp]
\centering
\diagram{windows}
\caption{Dependence across the three cross-sections. Bars depict trailing annual market windows schematically; fundamentals follow their recorded issuer reporting dates, not these uniform windows.}\label{fig:windows}
\end{figure}


Across the three dates, the dataset contains 123 instrument-date records and 43 distinct tickers. Forty instruments continue between each adjacent pair, and 39 are present at all three dates. Aspen is observed only in June 2025, Mondi at the first two dates, Glencore at the last two and Pan African Resources only in June 2026. The full three-date membership diagram therefore contains 39 instruments, while the clustering analysis retains all 41 instruments at each date. Because firms recur, annual windows overlap and some issuer accounts are used more than once, the 123 records do not represent 123 independent companies.


\subsection{Financial feature representation}

The clustering input contains eight financial and market variables, summarised in Table~\ref{tab:2}. Returns, beta, volatility, dividend yield and earnings yield describe performance, risk and valuation. Log market capitalisation, log revenue relative to market capitalisation and log share turnover describe scale and trading activity. This representation was developed after exploratory analysis to reduce skewness and the effect of share subdivisions. It extends the earlier study but uses a different sample and feature definition from the original eleven-variable SACJ analysis. The archived ten-variable specification is retained as a separate reference.


\begin{table}[htbp]
\centering\caption{Variables used by the current study model. Sector, company and ticker identify or describe observations and are excluded from fitting.}\label{tab:2}
\footnotesize
\begin{tabularx}{\linewidth}{@{}>{\hsize=0.560\hsize\linewidth=\hsize\raggedright\arraybackslash}X >{\hsize=1.440\hsize\linewidth=\hsize\raggedright\arraybackslash}X@{}}
\toprule
Feature & Construction and interpretation \\
\midrule
One-year return & Trailing split-adjusted price change; cash dividends excluded. \\
Beta & Slope relative to daily J203 market returns over the trailing year. \\
Annualised volatility & Daily-return sample standard deviation multiplied by the square root of 252. \\
Dividend yield & Trailing dividends relative to the security close, under the recorded currency convention. \\
Earnings yield & EPS / close; clipped at June 2025 first and ninety-ninth percentiles. \\
Log market capitalisation & Base-10 logarithm of rand line-specific market capitalisation. \\
Log revenue / capitalisation & Base-10 logarithm of revenue divided by market value; group totals for Investec. \\
Log share turnover & Base-10 logarithm of corrected average daily JSE volume / line-specific shares. \\
\bottomrule
\end{tabularx}
\end{table}


The treatment of share subdivisions is particularly important for comparisons over time. For a subdivision factor $q$, price and EPS are divided by $q$, while shares and consistently measured share volume are multiplied by $q$. Earnings yield, market capitalisation, revenue relative to market capitalisation and share turnover are unchanged when these quantities are measured on a consistent basis. However, a trailing-volume series can contain both old and new share units. Such a series must be corrected before turnover is calculated.


\subsection{Corrections and missing observations}

The original source values were preserved, and corrections were applied through separate reference tables. Naspers began trading on a five-for-one share basis on 1 October 2025, with the register updated on 6 October. Its trailing average volume contained mixed share units. The December 2025 value was corrected from 800,935 to 2,140,029 post-subdivision shares, and the June 2026 value from 1,502,524 to 1,847,639. These corrections place volume on the same share basis as the corresponding shares outstanding.


The Investec observations required a separate adjustment because both rows originally carried the combined share count. Issuer disclosures from June 2024 and June 2026 were used to allocate the reconstructed group total between the two listed lines. The December 2025 allocation uses the midpoint of the disclosed fractions. Each line then has its own market capitalisation and turnover. Revenue relative to market capitalisation uses group totals because each line represents a claim on the dual-listed group. The reconstructed total and the interpolated December allocation remain sources of measurement uncertainty.


Missing observations were handled without removing instruments from the sample. Reinet has no usable raw revenue or EPS at any date, while Anglo American lacks usable volume at the first two dates. Model values were imputed using June 2025 medians. Where a derived feature depends on an imputed input, its own June 2025 median was used. Earnings yield was clipped at the June 2025 first and ninety-ninth percentiles. These imputation and clipping values were held constant at the later dates.


Market measures were reconstructed from vendor daily histories. Dated issuer releases support revenue and diluted EPS for 120 of the 123 records, with dated currency conversions, while historical shares and market capitalisation remain reconstructed. Sector labels come from a legacy eight-category directory classification and were excluded from model fitting and selection. JSE-only volume may understate total trading activity for firms with major overseas listings, and revenue definitions differ across sectors. These data limitations are considered when interpreting the results.
